\documentclass[11pt]{article}
\usepackage[margin=1in]{geometry}
\usepackage{amsmath,amssymb}
\title{Disproof of Conjecture 00000007684}
\author{earthking11}
\date{6 October 2026}
\begin{document}
\maketitle

Choose $q=1/2$, which lies in the usual $q$-difference domain
$0<|q|<1$.  At the conjecture's listed scalar spectral value $\lambda=0$,
the eigenvalue equation is
\[
 D_q f(z)=0,
 \qquad\text{equivalently}\qquad f(qz)=f(z).
\]
Let the analytic function have power series $f(z)=\sum_{n\ge0}c_nz^n$.
Uniqueness of power-series coefficients gives
\[
 (q^n-1)c_n=0\quad(n\ge0).
\]
For every $n\ge1$, $q^n\ne1$, so $c_n=0$.  Therefore $f$ is constant.
The boundary condition now says
\[
 0=f(1)=c_0,
\]
so $f$ is the zero function.  The scalar solution space at $\lambda=0$ is
therefore the zero vector space and has dimension $0$, contradicting the
asserted value $d(0)=1$.

The Lean file works with coefficient sequences.  For $q=1/2$, the normalized
positive-degree coefficient equation is $c_n=0$ for $n>0$; the boundary
equation is $c_0=0$.  Lean proves that every sequence satisfying both is
identically zero, and that the zero sequence satisfies them.  This exactly
formalizes the algebraic coefficient step after the standard uniqueness
theorem for analytic power series.
\end{document}
